\documentclass[11pt,a4paper]{letter}
\usepackage[margin=1in]{geometry}
\usepackage{hyperref}
\signature{Ittipon Fongkaew \\ Corresponding author}
\address{School of Physics, Institute of Science \\ Suranaree University of Technology \\ 111 University Avenue \\ Nakhon Ratchasima 30000, Thailand \\ \texttt{ittipon@sut.ac.th}}
\begin{document}

\begin{letter}{Editor-in-Chief \\ \emph{Applied Energy} \\ Elsevier}

\opening{Dear Editor,}

We are pleased to resubmit, as a revised manuscript, our paper \emph{``A Probabilistic 48-Step Forecasting Benchmark for the Japan Electric Power Exchange: Conformal, Distributional, and Quantile-Averaging Methods Across Day-Ahead and Imbalance Markets''} for consideration as a research article in \emph{Applied Energy}. A detailed point-by-point response to each reviewer comment is attached as the document \texttt{RESPONSE\_TO\_REVIEWERS.pdf}.

The revision addresses all critical, major, and minor concerns raised in the first-round Editorial Decision Letter. Beyond those items, in response to reviewer interest in modern and zero-shot forecasting architectures we have substantially expanded the model set under the same Lago-protocol rolling-window evaluation. The revised manuscript benchmarks ten forecasters instead of the original six: the original LEAR, DDNN-Normal, DDNN-Johnson-SU, N-HiTS, QRA-rolling, and LEAR-plus-adaptive-conformal pipeline; plus the modern Transformer architectures PatchTST and iTransformer trained from scratch; plus the two pretrained zero-shot foundation forecasters Chronos-Bolt-small and TimesFM-1.0-200M. The MCS, block-length-sensitivity, point-accuracy, and probabilistic-accuracy results have all been re-run on the expanded model set.

The headline of the revised paper has substantially changed as a result. We summarise the new contributions below:

\begin{enumerate}
\item \textbf{The pretrained zero-shot foundation model Chronos-Bolt-small, evaluated with frozen weights and no JEPX-specific fine-tuning, is in the Hansen-Lunde-Nason Model Confidence Set on every (area, market) pair} at both \(\alpha = 0.10\) and \(\alpha = 0.25\), and is the sole MCS singleton on Tokyo IM. The result is robust to all three stationary-bootstrap mean block lengths we tested (24, 336 weekly, 672 biweekly half-hours). This is, to our knowledge, the first published instance of a zero-shot foundation forecaster statistically matching the strongest supervised method on a major spot electricity market.
\item \textbf{The modern Transformer architecture iTransformer substantially outperforms the Lago-tradition supervised models} (LEAR, DDNN-Normal, DDNN-Johnson-SU, N-HiTS) on every (area, market) pair, with rRMSE 0.72--0.76 versus 0.77--1.21. iTransformer is in the MCS on three of four pairs alongside Chronos-Bolt; the original supervised winners (N-HiTS on DA, LEAR on IM) are eliminated from the MCS once the new candidates enter the set.
\item \textbf{Foundation-model success on JEPX is not generic.} The second pretrained zero-shot foundation model we evaluated, TimesFM-1.0-200M, performs poorly even after physical-range clipping of its outputs (rRMSE 5--10 across cells), driven by occasional extreme tokenization outliers in its autoregressive decoder. We document this failure mode in the manuscript as a methodological caution for the EPF community.
\item \textbf{The original Lago-tradition supervised-only 2$\times$2 cross-market identification (N-HiTS on DA, LEAR on IM)} is preserved in the manuscript as a sub-finding, with the explicit observation that it characterises the restricted Lago-tradition candidate set but does not survive when modern Transformer or foundation-model candidates are added.
\item \textbf{All numerical inconsistencies flagged by the Editor and the Devil's Advocate reviewer have been corrected,} including the §1 / Table 4 contradiction on the Tokyo DA MCS finding, the missing values in Table 1, the implicit Winkler-claim reversal in the original Abstract, the DM-HAC lag selection, the CRPS-comparability across closed-form and empirical-CDF estimators, the conformal calibration-vs-scoring window separation, and the orphan dependency declarations.
\end{enumerate}

The four originally-flagged critical issues have all been resolved. We have also added (i) a Highlights file, (ii) an Ethics statement, (iii) the necessary citations for the new model architectures, (iv) a methodological note in the Discussion on look-ahead bias in P\&L back-tests (foregrounding the v1e preprint retraction that was disclosed in the original cover letter), and (v) an updated Data and Code Availability statement reflecting the expanded reproducibility package.

We have re-verified every numerical value in the revised manuscript against the released metric files, and the manuscript is fully reproducible from the released code and per-issuance prediction matrices. The reproducibility package is at the redacted-author GitHub mirror and on a reserved Zenodo deposit, both of which will be finalised on acceptance.

We believe the revised manuscript represents a substantially stronger contribution than the original submission, and that the Chronos-Bolt zero-shot result in particular is of broad interest to the \emph{Applied Energy} readership working on spot-market forecasting and on the recent foundation-model wave.

Thank you for handling this revised submission. We look forward to your response.

\closing{Sincerely,}

\end{letter}
\end{document}
